%!tex

\input tikz
\usetikzlibrary{arrows.meta, math}

\input ../../../miscelania/macros.tex

\xiiptrm

\null\vfill

\line{%
  \hfil\xivbx
  Experimento: segunda lei de Newton
  \hfil
}
\vskip .5cm

\fig{%
  \tikzpicture[]
     \coordinate (C1) at (2.2,  6.945)            ;
     \coordinate (C2) at (4.4,  9.647)            ;
     \coordinate (C3) at (6.6, 11.761)            ;
     \coordinate (C4) at (8.8, 12.901)            ;
     \coordinate (D1) at (2.2, 2.635274371087693);
     \coordinate (D2) at (4.4, 3.106018107617611);
     \coordinate (D3) at (6.6, 3.429476459402069);
     \coordinate (D4) at (8.8, 3.591769500816717);
     \draw (D1) -- (D4);
     \foreach \i in {1,...,4} {
       \draw (C\i) -- (D\i);
     };
     \fill[red  ] (C1) circle [radius=1.5pt];
     \fill[red  ] (C2) circle [radius=1.5pt];
     \fill[red  ] (C3) circle [radius=1.5pt];
     \fill[red  ] (C4) circle [radius=1.5pt];
     \fill[green] (D1) circle [radius=1.5pt];
     \fill[green] (D2) circle [radius=1.5pt];
     \fill[green] (D3) circle [radius=1.5pt];
     \fill[green] (D4) circle [radius=1.5pt];
  \endtikzpicture
}{%
  gr\'aficos.
}

\vfill\null

\bye
